\section*{الفصل \ref{ch:b1:e-ph-diagrams} --- مخططات الجهد--pH}

\begin{solution}{exo:b1:e-ph-diagrams:1}
المنغنيز: \babelsublr{\ce{Mn} (0) $<$ \ce{Mn^2+}, \ce{Mn(OH)2} ($+$II) $<$ \ce{MnO2}
($+$IV) $<$ \ce{MnO4-} ($+$VII)}. الكلور: \babelsublr{\ce{Cl-} ($-$I) $<$ \ce{Cl2} (0)
$<$ \ce{HClO}, \ce{ClO-} ($+$I) $<$ \ce{ClO3-} ($+$V)}.
\end{solution}

\begin{solution}{exo:b1:e-ph-diagrams:2}
$-0.059 \times 3/1 = -0.177$؛ $+0.059 \times 2/2 = +0.059$؛
$-0.059 \times 8/5 = -0.094$؛ $-0.059 \times 2/2 = -0.059$ V لكل وحدة pH.
وفي مزدوجة النحاس تتحرر البروتونات في جهة الشكل
\omterm{def:b1:nernst:couple}{المختزل} ($m = -2$): فيرتفع الجهد مع pH.
\end{solution}

\begin{solution}{exo:b1:e-ph-diagrams:3}
تمر $\ce{H+}/\ce{H2}$ بالقيمة 0~V عند pH 0 ولا تبلغ أبدًا 0.50~V من أجل
$\mathrm{pH} \geqslant 0$. وتبلغ $\ce{O2}/\ce{H2O}$ القيمة 0.50~V عند
$\mathrm{pH} = 0.73/0.059 = 12.4$ ولا تبلغ 0~V إلا عند pH 20.8، خارج
السلّم. وعرض الشريط \qty{1.23}{V} عند كل pH، ومنها pH 7
(من $-0.41$ إلى $+0.82$~V).
\end{solution}

\begin{solution}{exo:b1:e-ph-diagrams:4}
(1، 1.0~V): \ce{Fe^3+}. (4، 0): المستقيم $\ce{Fe(OH)3}/\ce{Fe^2+}$ عند
$1.09 - 0.71 = 0.38$~V والنقطة تحته، وفوق $-0.47$~V:
\ce{Fe^2+}. (10، $-0.4$~V): بين $-0.07 - 0.59 = -0.66$ و$0.28 -
0.59 = -0.31$~V: \ce{Fe(OH)2}. (10، 0.5~V): \ce{Fe(OH)3}.
\end{solution}

\begin{solution}{exo:b1:e-ph-diagrams:5}
$\ce{Fe^2+}/\ce{Fe}$: $-0.41 + 0.030 \times (-6) = -0.59$~V.
$\ce{Fe^3+}/\ce{Fe(OH)3}$: $14.00 - \frac13(38.55 - 6) = 3.15$. تتسع مناطق
الأيونات المذابة (التآكل): فتمتد \ce{Fe^2+} نزولًا إلى $-0.59$~V
و\ce{Fe^3+} إلى pH 3.15؛ ويخسر الفلز والهيدروكسيدان من مساحتهما.
\end{solution}

\begin{solution}{exo:b1:e-ph-diagrams:6}
\ce{Fe(OH)3 + 3H+ + e- -> Fe^2+ + 3H2O}: $E = E^\circ + 0.059\log\frac{[\ce{H+}]^3}{[\ce{Fe^2+}]}
= (E^\circ - 0.059\log c) - 0.177\,\mathrm{pH}$. وعند pH 1.82 يلتقي المستقيم
$\ce{Fe^3+}/\ce{Fe^2+}$ عند 0.77~V، فيكون الثابت $0.77 + 0.177 \times 1.82
= 1.09$~V.
\end{solution}

\begin{solution}{exo:b1:e-ph-diagrams:7}
$E^\circ(\ce{Cu+}/\ce{Cu}) = 0.52 > E^\circ(\ce{Cu^2+}/\ce{Cu+}) = 0.16$: فحسب
\cref{prop:b1:e-ph-diagrams:disproportionation-criterion} يخضع \ce{Cu+}
\omterm{def:b1:e-ph-diagrams:disproportionation}{لعدم التناسب}. والحد الوحيد: $E = 0.34 + 0.030\log 0.010 =
0.28$~V.
\end{solution}

\begin{solution}{exo:b1:e-ph-diagrams:8}
عند pH 0 تتداخل منطقة النحاس (تحت 0.28~V) مع منطقة الماء
وتقع فوق مستقيم الهيدروجين: فلا تفاعل مع حمض الكلوريدريك
بمعزل عن الهواء. وحمض النتريك: $\ce{NO3-}/\ce{NO}$ عند 0.96~V يقع فوق
منطقة النحاس، فيتأكسد النحاس:
\ce{3Cu + 2NO3- + 8H+ -> 3Cu^2+ + 2NO + 4H2O}.
% ledger: dfg:NO3-_ao, dfg:NO_g (E0 computed in figdata/bachelor-1/standard-potentials.py)
\end{solution}

\begin{solution}{exo:b1:e-ph-diagrams:9}
عند pH 3 يقع المستقيم $\ce{O2}/\ce{H2O}$ عند $1.23 - 0.18 = 1.05$~V، فوق
منطقة \ce{Fe^2+} كلها: فيؤكسد الأكسجين الحديد(II). وفوق pH 1.82 يترسب
الحديد(III) المتكوّن على هيئة الهيدروكسيد الأصفر البني:
\ce{4Fe^2+ + O2 + 10H2O -> 4Fe(OH)3 + 8H+}.
\end{solution}

\begin{solution}{exo:b1:e-ph-diagrams:10}
الحدان الرأسيان: $14.00 - \frac12(16.38 - 2) = 6.81$؛ و، انطلاقًا من
\ce{Zn(OH)2 + 2OH- <=> Zn(OH)4^2-} ($K = 10^{-1.93}$) مع
$[\ce{Zn(OH)4^2-}] = 10^{-2}$، $[\ce{OH-}]^2 = 10^{-0.07}$، pH 13.97.
$\ce{Zn^2+}/\ce{Zn}$: $-0.76 + 0.030 \times (-2) = -0.82$~V. $\ce{Zn(OH)2}/\ce{Zn}$:
الميل $-0.059$، ويمر بالنقطة $(6.81, -0.82)$: $E = -0.42 - 0.059\,\mathrm{pH}$.
\ce{Zn(OH)4^2- + 4H+ + 2e- -> Zn + 4H2O}: الميل $-0.118$، ويمر بالنقطة
$(13.97, -1.24)$: $E = 0.41 - 0.118\,\mathrm{pH}$.
\end{solution}

\begin{solution}{exo:b1:e-ph-diagrams:11}
تقع منطقة الحديد تحت مستقيم الهيدروجين عند كل pH. عند pH 2: يتأكسد
الحديد إلى \ce{Fe^2+} مع انطلاق الهيدروجين (تآكل). وعند pH 7:
يبقى التفاعل ممكنًا، لكن مستقيم الهيدروجين ($-0.41$~V)
قريب من $\ce{Fe^2+}/\ce{Fe}$ والتفاعل بطيء جدًا. وعند pH 13: يتأكسد الحديد
إلى \ce{Fe(OH)2}، وهي \omterm{def:b1:e-ph-diagrams:corrosion-domains}{منطقة تخميل}: فتتكوّن كسوة ويحتفظ
المسمار ببريقه. ومع الأكسجين المذاب، يؤكسد المستقيم $\ce{O2}/\ce{H2O}$،
الواقع بعيدًا فوق، الحديدَ إلى حديد(III) عند كل pH: إنه الصدأ، \ce{Fe(OH)3}
وأكاسيده المنزوعة الماء.
\end{solution}

\begin{solution}{exo:b1:e-ph-diagrams:12}
منطقة الزنك هي الأدنى: فهو يتأكسد أولًا حين يتماس مع الفولاذ،
\ce{2Zn + O2 + 2H2O -> 2Zn(OH)2} عند pH 8، والإلكترونات التي
يحررها تختزل الأكسجين على سطح الفولاذ، الذي يبقى فلزيًا.
والمغنيسيوم ($E^\circ = -2.36$~V) أدنى أيضًا وينفع كذلك؛ أما النحاس
(0.34~V) فيقع فوق الحديد: فإذا وُصل الحديد بالنحاس كان الحديد هو الفلز
المتأكسد، وأسرع مما لو كان وحده.
\end{solution}

\begin{solution}{pb:b1:e-ph-diagrams:1}
\textbf{1.} \babelsublr{$-$I}، 0، \babelsublr{$+$I}، \babelsublr{$+$I}.
\textbf{2.} \ce{Cl-} في الأسفل، ثم \ce{Cl2}، ثم \ce{HClO} و\ce{ClO-}
في الأعلى.
\textbf{3.} \ce{HClO} و\ce{ClO-}، وهما مزدوجة حمض--قاعدة.
\textbf{4.} \ce{Cl2 + 2e- -> 2Cl-}.
\textbf{5.} \ce{2HClO + 2H+ + 2e- -> Cl2 + 2H2O}.
\textbf{6.} \ce{ClO- + H2O + 2e- -> Cl- + 2OH-}.
\textbf{7.} \ce{HClO} تحت pH 7.55، و\ce{ClO-} فوقه.
\textbf{8.} لا يُتبادل أي إلكترون؛ فالحد محدد بالثابت $K_a$ وحده.
\textbf{9.} $1/(1 + 10^{7.4 - 7.55}) = 0.59$: أي \qty{59}{\percent}.
\textbf{10.} $1/(1 + 10^{0.45}) = 0.26$: فثلاثة أرباع الكلور الفعال
ستكون على هيئة \ce{ClO-}، المطهِّر الأضعف.
\textbf{11.} \ce{ClO-}.
\textbf{12.} $E = 1.396 + 0.0296\log\frac{[\ce{Cl2}]}{[\ce{Cl-}]^2} = 1.396
+ 0.0296\log\frac{0.005}{10^{-4}} = 1.396 + 0.050 = 1.446$~V؛ ولا \ce{H+}
في \omterm{def:b1:nernst:couple}{نصف المعادلة}.
\textbf{13.} $E = 1.594 + 0.0296\log\frac{[\ce{HClO}]^2[\ce{H+}]^2}{[\ce{Cl2}]}
= 1.594 + 0.0296\log\frac{10^{-4}}{0.005} - 0.0592\,\mathrm{pH} = 1.594 -
0.050 - 0.0592\,\mathrm{pH}$.
\textbf{14.} بالمساواة، $0.0592\,\mathrm{pH} = 1.594 - 1.396 - 2 \times 0.0503
= 0.0974$، $\mathrm{pH} = 1.646$.
\textbf{15.} \ce{HClO + H+ + 2e- -> Cl- + H2O}: بروتون واحد لإلكترونين،
والميل $-0.0592/2 = -0.0296$.
\textbf{16.} بجمع نصفي المعادلتين (إلكترون واحد لكل ذرة كلور
في كل منهما): $2E^\circ = 1.594 + 1.396$، $E^\circ(\ce{HClO}/\ce{Cl-}) =
1.495$~V؛ والحد $E = 1.495 - 0.0296\,\mathrm{pH}$.
\textbf{17.} \ce{ClO- + 2H+ + 2e- -> Cl- + H2O}: $E = 1.718 - 0.0592\,\mathrm{pH}$
(الثابت مختار للاتصال: $1.495 - 0.0296 \times 7.55 = 1.272 = 1.718
- 0.0592 \times 7.55$).
\textbf{18.} \ce{Cl-} تحت المستقيمات؛ و\ce{Cl2} في مثلث صغير عند
$\mathrm{pH} < 1.65$ فوق 1.446~V؛ و\ce{HClO} فوق المستقيم ذي الميل
$-0.0296$ حتى pH 7.55، و\ce{ClO-} بعدها. والمستقيم $\ce{O2}/\ce{H2O}$،
من 1.23 إلى 0.40~V، يقع تحت كل هذه الحدود (وللمخطط
المحسوب ببرنامج الفصل الشكل نفسه).
\textbf{19.} لا: منطقته تقع فوق مستقيم الأكسجين، فيمكنه أن
يؤكسد الماء، لكن التفاعل بطيء. ويفقد المسبح كلوره أساسًا
لأن الكلور يؤكسد ما يجلبه السبّاحون والهواء، وبفعل
ضوء الشمس؛ فيجب تعويضه.
\textbf{20.} ليست للنوع \ce{Cl2} منطقة هناك: فهو يخضع \omterm{def:b1:e-ph-diagrams:disproportionation}{لعدم التناسب}،
\ce{Cl2 + H2O -> HClO + Cl- + H+}؛ وفي القاعدة
\ce{Cl2 + 2OH- -> ClO- + Cl- + H2O}.
\textbf{21.} بإمرار الكلور في محلول هيدروكسيد الصوديوم (تفاعل
السؤال 20).
\textbf{22.} \ce{HClO + Cl- + H+ -> Cl2 + H2O}، وهو تفاعل \omterm{def:b1:e-ph-diagrams:disproportionation}{تناسب}.
\textbf{23.} تحت pH 1.6 يتناسب الهيبوكلوريت والكلوريد: فيتكوّن الكلور،
وما زاد على ذوبانيته ينطلق غازًا سامًا.
\textbf{24.} عند $c = \qty{0.010}{mol/L}$ يقع pH 4 خارج منطقة
\ce{Cl2}. لكن الحد يتعلق بالتركيز $c$: فإعادة الأسئلة 12--14 من أجل
$c$ عام تعطي $\mathrm{pH} = 3.34 + \log(2c)$، الذي يبلغ 3.6 من أجل
$c = \qty{1}{mol/L}$، وهي رتبة تركيز ماء جافيل. فخلط
المنتجات المركزة قرب pH 4 ليس مأمونًا هو أيضًا.
\textbf{25.} \textbf{pH 1.6}: فوقه، عند $c = \qty{0.010}{mol/L}$،
يخضع الكلور المذاب لعدم التناسب معطيًا حمض الهيبوكلوروز والكلوريد.
\end{solution}
